\theory{Singularity theory}{How do spaces, maps, and physical patterns change when a smooth or regular situation degenerates?}{Singularity theory studies points where manifold structure fails, critical points of maps, and sudden changes under parameter variation. It classifies stable local models and studies how singularities arise by projection, symmetry, or degeneration.}{Algebraic geometry, differential geometry, wave fronts, optics, mechanics, catastrophe theory, bifurcation theory, and intersection cohomology.}{Jacobian rank, local normal forms, and equivalence under coordinate changes classify where smooth behavior breaks down.}

\paragraph{The idea and history}
A string is a one-dimensional manifold until it crosses itself. A crossing is a double point; a touching point can disappear after a tiny movement and is unstable. Singularity theory asks which irregularities persist under small changes and which are artifacts of a special drawing \cite{singularity_ref1}.

Vladimir Arnold described the subject as the study of how objects depend on parameters when a small change causes a sudden qualitative change. Whitney studied stable singularities and critical points; Rene Thom developed catastrophe theory; Arnold connected the subject with symplectic geometry and mechanics \cite{singularity_ref2}.

\end{historylist}
\section*{3. Theory}
\subsection*{Core theory}
\paragraph{Algebraic and geometric examples}
The curve $y^2=x^2+x^3$ has a double point at the origin, while $y^2=x^3$ has a cusp. A singular point of a variety is one where the local equations do not define a smooth manifold. The Jacobian matrix detects this: the rank is too small at a singular point.

Singularities also arise by projecting a smooth object into lower dimensions. Caustics, such as bright patterns at the bottom of a swimming pool, occur where many light rays focus. Symmetry can create orbifolds with corners or cone-like points. A smooth manifold can degenerate into a singular space as parameters change \cite{singularity_ref3}.

\paragraph{Stability, germs, and catastrophes}
To classify a local singularity, one studies a germ, the behavior near a point, up to changes of coordinates. Taylor expansions and jets record derivatives up to a chosen order. Group actions identify equivalent local forms.

Catastrophe theory studies stable critical-point changes in families of smooth functions. Fold, cusp, swallowtail, and butterfly are standard elementary models. The models do not claim that every natural event is one of these shapes; they describe local mechanisms by which equilibria appear, merge, or disappear \cite{singularity_ref1}.

\paragraph{Resolution and duality}
Hironaka's resolution-of-singularities theorem in characteristic zero says that singularities of algebraic varieties can be replaced by a smooth variety through a controlled sequence of transformations. The operation is not simply erasing a bad point; it records the exceptional geometry introduced during the resolution.

Singular spaces may fail Poincare duality. Intersection cohomology restores a form of duality by using the stratification of the singular space. Perverse sheaves arise from this framework and connect singularity theory with homological algebra. Monodromy studies how solutions or covering spaces behave around singularities of differential equations \cite{singularity_ref4}.

\section*{4. Important theorems and results}
\begin{enumerate}[leftmargin=*,itemsep=0.35em]
\item \textbf{Jacobian criterion.} Under standard hypotheses, a point of an algebraic variety is smooth exactly when the Jacobian has the expected rank.

\item \textbf{Whitney stability theory.} Stable singularities are those whose qualitative form persists under small perturbations.

\item \textbf{Morse lemma.} Near a nondegenerate critical point, a smooth function is equivalent to a quadratic form.

\item \textbf{Hironaka resolution theorem.} Algebraic varieties over characteristic zero admit resolutions of singularities.

\item \textbf{Intersection-cohomology duality.} Intersection cohomology restores a Poincare-type duality for broad classes of singular spaces.

\item \textbf{Classification of elementary catastrophes.} Low-dimensional stable critical changes reduce locally to a short list of normal forms.
\end{enumerate}
\noindent\emph{Source for these results:} \cite{singularity_ref1}.

\theoryapplications
\theoryconnections{singularity_theory}{%
\item Calculus identifies where a map loses rank or a surface ceases to be smooth.
\item Differential and algebraic geometry classify the local equations, topology tests whether the local type survives small changes, and Lie group actions remove coordinate choices that do not change the singularity.
\item Homological tools then measure how cycles, intersections, or vanishing pieces appear when the singular point is resolved or deformed \cite{singularity_ref1,singularity_ref2}.
}{%
\item Singularity theory supplies normal forms for catastrophe and bifurcation models, so a complicated equation can be replaced locally by a short list of stable patterns.
\item In optics, the singular set of a ray map produces caustics and wave fronts; in algebraic geometry, resolutions and intersection cohomology recover information hidden by a singular point.
\item The same classification therefore links visible folds and cusps with algebraic invariants \cite{singularity_ref1,singularity_ref3}.
}
\section*{Glossary}
\begin{description}[leftmargin=*,style=nextline,itemsep=0.3em]
\item[\textbf{Singularity.}] A point where a mathematical object (such as a curve, surface, or map) fails to be smooth or regular; for example, a cusp or crossing point.
\item[\textbf{Germ.}] The local behavior of a function or map near a specific point, considered up to changes of coordinates.
\item[\textbf{Caustic.}] A visible pattern formed by the focusing or envelope of light rays; it appears as a bright curve or surface where many rays meet, creating a singularity in the projection.
\item[\textbf{Resolution of singularities.}] A process that replaces a singular variety with a smooth one while keeping track of the exceptional sets introduced, effectively "smoothing out" the singularities.
\item[\textbf{Intersection cohomology.}] A tool that restores Poincaré duality for singular spaces by restricting chains near singular strata, providing a way to study the topology of singular objects.
\item[\textbf{Morse lemma.}] A result stating that near a nondegenerate critical point, a smooth function can be expressed as a quadratic form in appropriate local coordinates.
\item[\textbf{Catastrophe theory.}] The study of how small changes in parameters can cause sudden qualitative changes in the behavior of a system, especially in families of smooth functions.
\item[\textbf{Stratification.}] A decomposition of a singular space into smooth manifold pieces arranged so each piece sits cleanly on top of lower-dimensional ones.
\item[\textbf{Monodromy.}] The action on solutions obtained by analytic continuation around a singular point.
\item[\textbf{Normal form.}] A canonical representative of a singularity up to coordinate changes, used to classify qualitatively distinct local behaviors.
\item[\textbf{Jets.}] The Taylor expansion of a map up to a chosen finite order, recording derivatives that distinguish nearby maps.
\end{description}
\section*{7. Sample Questions}
\emph{Exercise reference:} \cite{singularity_ref1}. The cited edition is the source for the exercise set; consult its Exercises section for the official statement and numbering.
\begin{enumerate}[leftmargin=*,itemsep=0.9em]
\item \textbf{Exercise 1.} Determine whether the origin is a singular point of the curve $y^2 = x^2 + x^3$. \emph{Solution.} Let $F(x,y) = y^2 - x^2 - x^3$. The gradient is $\nabla F = (-2x - 3x^2,\; 2y)$. At the origin, $\nabla F(0,0) = (0,0)$, so the Jacobian criterion identifies the origin as a singular point. Geometrically, the curve has a node (self-crossing) at the origin.

\item \textbf{Exercise 2.} Determine whether the origin is singular on $y^2 = x^2 + x^4$ and on $y^2 = x^4 + x^6$. \emph{Solution.} For $F_1(x,y) = y^2 - x^2 - x^4$: $\nabla F_1(0,0) = (0,0)$, so the origin is singular (a node). For $F_2(x,y) = y^2 - x^4 - x^6$: $\nabla F_2(0,0) = (0,0)$ as well, so this is also singular. However, after blowing up once, $F_1$ becomes $v^2 = 1 + u^2$ (smooth), while $F_2$ becomes $v^2 = u^2 + u^4$ (still singular). Thus $F_1$ has a simpler singularity requiring one blow-up, while $F_2$ requires two.

\item \textbf{Exercise 3.} By the Morse lemma, near a nondegenerate critical point a smooth function $f\colon\mathbb{R}^2\to\mathbb{R}$ is equivalent to a quadratic form. Compute the Hessian of $f(x,y) = x^2 + 3y^2 - 2xy$ at the origin, classify the critical point, and state the normal form. \emph{Solution.} The partial derivatives are $f_x = 2x - 2y$ and $f_y = 6y - 2x$, which both vanish at the origin. The Hessian matrix is $H = \begin{pmatrix} 2 & -2 \\ -2 & 6 \end{pmatrix}$. Its determinant is $2 \cdot 6 - (-2)^2 = 8 > 0$ and $f_{xx} = 2 > 0$, so the origin is a nondegenerate minimum. By the Morse lemma, $f$ is locally equivalent to $u^2 + v^2$ after a suitable change of coordinates.

\item \textbf{Exercise 4.} Verify that the curve $y^2 = x^5$ has a cusp at the origin by computing its tangent cone and comparing it with the node on $y^2 = x^2(x+1)$. \emph{Solution.} For $y^2 = x^5$, the lowest-degree terms in $F(x,y) = y^2 - x^5$ are $y^2$, so the tangent cone is the double line $y^2 = 0$ (the $x$-axis counted twice). This is a cusp. For $y^2 = x^2(x+1) = x^3 + x^2$, the lowest-degree terms are $y^2 - x^2 = (y-x)(y+x)$, giving two distinct tangent lines $y = x$ and $y = -x$. This is a node (transverse self-crossing).

\item \textbf{Exercise 5.} Compute the Euler characteristic of the blow-up of $\mathbb{CP}^2$ at a single point. \emph{Solution.} The Euler characteristic of $\mathbb{CP}^2$ is $\chi(\mathbb{CP}^2) = 3$ (since it has a cell decomposition with cells in dimensions $0, 2, 4$). Blowing up a point replaces it with an exceptional $\mathbb{CP}^1$, contributing $\chi(\mathbb{CP}^1) = 2$ instead of the single point's contribution of $\chi(\text{pt}) = 1$. Thus $\chi(\text{blow-up}) = 3 - 1 + 2 = 4$. Equivalently, the blow-up is $\mathbb{CP}^2 \# \overline{\mathbb{CP}}^2$, and $\chi = 3 + 3 - 2 = 4$.

\end{enumerate}
\section*{8. References}
\begin{thebibliography}{9}
\bibitem{singularity_ref1} V. I. Arnold, \emph{Singularities of Differentiable Maps}, Birkhauser, 1982.
\bibitem{singularity_ref2} V. I. Arnold, \emph{Catastrophe Theory}, 2nd ed., Springer, 1986.
\bibitem{singularity_ref3} E. Brieskorn and H. Knorrer, \emph{Plane Algebraic Curves}, Birkhauser, 1981.
\bibitem{singularity_ref4} M. Goresky and R. MacPherson, \emph{Stratified Morse Theory}, Springer, 1980.
\end{thebibliography}
